\section{Conditions on a deformation mechanism}\label{sec:conditions}

A candidate arising from this picture inherits the necessary conditions the
companion investigation has accumulated \cite{Selection}. We do not restate their
derivations; the table records which ones bear on a flow and how.

\begin{center}
\begin{tabular}{@{}p{0.30\textwidth}p{0.62\textwidth}@{}}
\toprule
Inherited condition & Bearing on a deformation mechanism\\
\midrule
Reflection invariance & The flow must commute with $x\mapsto-x$; free to check,
and the involution is the functional equation acting on the zeros.\\
Unboundedness relative to $L^2$ & $Q_L[\chi e^{iNx}]/\log N\to\norm\chi^2$, so no
bounded preparation matches. This is what excluded the closed benchmark;
Condition~\ref{cond:unbounded} restates it for a holonomy.\\
Spectral weight and channel tower & The archimedean multiplier is the smooth zero
density to four orders, with the exact tower $\{2n+\frac12\}$ and unit residues.
Condition~\ref{cond:marginal} reads this as a defect two-point function.\\
Criticality & The relevant contraction is saturated and the margin collapses
superexponentially in $L$; any mechanism supplying a margin uniform in $L$ proves
something false. See Section~\ref{sec:outlook}, where we ask whether the flow
formulation is the exception.\\
Jump structure & The non-pole part of the target is one L\'evy jump form with
atoms at $m\log p$. Equation \eqref{eq:atomic} realizes this on the connection.\\
\bottomrule
\end{tabular}
\end{center}

We add three conditions specific to the deformation picture. They are
\emph{proposed}: each is elementary, none has been tested against a specific
theory, and the third carries an identification hypothesis that this program
deliberately does not make.

\begin{condition}[Unbounded connection, infinite-dimensional
representation]\label{cond:unbounded}
A holonomy valued in a compact group, in a finite-dimensional representation, is
unitary; its matrix elements and its character are bounded functions, so every
pairing assembled from them is bounded on $L^2(I_L)$ and is excluded by the
unboundedness rule. A deformation mechanism must therefore carry an unbounded
connection in an infinite-dimensional representation.
\end{condition}

This is the structural form of an exclusion the program already had. The closed
benchmark failed because its operator acted as $W=2\cos\theta$, with spectrum in
$[-2,2]$, so that its kernel was entire in the separation and had no atoms
\cite{Selection}. Condition~\ref{cond:unbounded} says that the boundedness was
not an accident of that benchmark: $2\cos\theta$ is the character of a holonomy
in the fundamental of $SU(2)$, and \emph{every} such character is bounded. Note
also that $V_{\omega,L}$ is not unitary --- it is a contraction, and the object
of interest is its defect \eqref{eq:defect} --- so the connection to be matched is
not anti-Hermitian either.

\begin{condition}[Reflection and the atoms]\label{cond:atoms}
The connection's kernel must be a continuous part plus atoms at the prime periods
$m\log p$ with masses $2(\log p)p^{-m/2}\cosh(\omega m\log p)$, and the whole must
be even in the separation after symmetrization. A connection that is smooth in
the separation cannot produce the atoms, and a connection whose atoms sit at
other positions is excluded without any positivity computation.
\end{condition}

\begin{condition}[Marginal short distance: effective dimension
$\frac12$]\label{cond:marginal}
Suppose the line's parameter is identified with the arithmetic coordinate $x$.
Since $\gammak(r)=\frac1{2r}+O(1)$ as $r\downarrow0$ by
\eqref{eq:gamma-kernel}, \eqref{eq:gamma-energy} gives
\[
 K[F]=\frac12\iint|F(x)-F(y)|^2\gammak(|x-y|)\dd x\dd y
 \sim\frac14\iint\frac{|F(x)-F(y)|^2}{|x-y|}\dd x\dd y
\]
near the diagonal: the Gagliardo seminorm at the borderline exponent
$1+2s=1$, that is $s=0$. An insertion on the line reproducing this must therefore
have two-point function $|x-y|^{-1}$, an effective defect dimension
$\Delta=\frac12$, with the coefficient fixed rather than free.
\end{condition}

Condition~\ref{cond:marginal} has one immediate consequence, and it excludes the
first thing one would try. The displacement operator of a line defect has
protected dimension $\Delta=2$, fixed by the Ward identity that defines it, and
the standard protected scalar insertions on a supersymmetric Wilson line have
$\Delta=1$ \cite{GiombiKomatsu,CHMS}. Neither is $\frac12$. \emph{So the
deformation whose generator matches is not the displacement operator}, which is
what a deformation of the contour would naively produce.

\begin{remark}[The identification hypothesis]\label{rem:identification}
Condition~\ref{cond:marginal} assumes the line's parameter is the arithmetic
coordinate. The program does not make that identification: the arithmetic
coordinate is not automatically a physical Euclidean time or contour parameter,
and an earlier calculation excluded one natural way of making it one
\cite{GaugeTransfer,InverseBulk}. Sections~\ref{sec:abelian} and
\ref{sec:endpoint} suggest a different assignment, in which $\omega$ is the
parameter along the line --- a cusp angle, by Remark~\ref{rem:angle} --- and $x$
indexes the representation the line carries. Under that assignment
Condition~\ref{cond:marginal} moves and must be rederived, while
Conditions~\ref{cond:unbounded} and \ref{cond:atoms} are unaffected. Deciding
between the two assignments is the first question a specific proposal has to
answer.
\end{remark}
